\section{Proof Machinery, Revision Law, Competition, and Compression}

The present release does not treat proof as a single binary flag. It carries an explicit proof machinery: a consistency layer, an empirical anchoring layer, a revision law, a minimality or irreducibility layer, an alternative-model competition layer, and a bounded compression benchmark.
This layer is what prevents the Core from becoming rhetorical rubber. The same machine-readable promotion law that constrains runtime status now constrains the monograph itself.

\subsection{Consistency and promotion law}
The foundational consistency discussion tracks the K-level doctrine rows as an evidence-boundary surface rather than as public proof by status token.
Positive outward science is permitted only for claims that are source-bound, proof-bound, support-consistent, and, when empirical, tied to observables, numerical packets, replay routes, and explicit falsifiers.

\subsection{Revision law}
The declared revision mutability is `DERIVED\_FIRST`: empirical bridges first, then parameterizations, then standalone domain packets and derived theorems, and only then -- under accumulated cross-domain failure -- the root kernel itself.
The current root-kernel revision status is `LOCKED\_PENDING\_ACCUMULATED\_CROSS\_DOMAIN\_FAILURE`.
This means that conflicts with reality do not license silent widening. They first force narrower repair at the empirical and derived layers. Any widening that moves the nonclaim boundary while preserving rhetoric is classified as illegal rubberization.

\subsection{Minimality and irreducibility}
The 1.3.3 minimality witness route records component-wise keep/drop witnesses for all nine canonical release-tuple components in the finite semantic suite, with source-inspected witness records for the corresponding typed witness schema.
The strongest honest meta-result remains component-wise independence for the declared semantic verdict suite, while frontier candidates beyond the declared tuple remain research obligations rather than part of the promoted T133-MIN theorem.

\subsection{Alternative-model competition}
The competition policy is `BOTH`: the Core must eventually compete both against same-claim-class meta-models and against domain baselines.
At present the matrix tracks `5` comparison rows, of which `4` are already pass-eligible.
The point of this layer is not branding but discipline: description length, support burden, and benchmark coverage must be compared under declared penalty rules rather than by prose.

\subsection{Compression as a bounded metric}
The release therefore treats compression as `BOUNDED\_METRIC` rather than as a central theorem claim.
The current benchmark exposes coverage per kernel statement `0.294118` and coverage per active numerical packet `1.0`.
These numbers are reported as bounded scientific metrics only. They do not by themselves promote the Core unless competition and empirical gates also pass.
